\section*{Chapitre \ref{ch:m2:market-making-games} --- Jeux de teneur de marché}

\begin{solution}{exo:m2:market-making-games:1}
Autour de l'espérance, 3.5~: par exemple «~3 à 4, 10 vers le haut~». La valeur est connue, donc personne
ne peut être mieux informé~; la largeur représente un petit profit pour le risque de roll et une
largeur à laquelle l'autre côté veut encore échanger.
\end{solution}

\begin{solution}{exo:m2:market-making-games:2}
Relever le marché~: l'acheteur en sait peut-être davantage, et il a au moins jugé 400 bon marché. Déplacer
le centre vers le haut et peut-être élargir, par exemple en cotant 300 à 500, et réfléchir à la raison pour laquelle elle a acheté~:
si elle a pu compter les fenêtres, il faut beaucoup le déplacer.
\end{solution}

\begin{solution}{exo:m2:market-making-games:3}
Un trader informé passe lorsque la somme est comprise entre 30 et 40, et un trader non informé passe
avec la même probabilité quelle que soit la somme~; la loi a priori est symétrique autour de 35, et il en va de même
l'ensemble des sommes qui produisent un passage, donc l'espérance ne bouge pas. Plus tard, lorsque les cotations sont
décentrées, un passage la déplace.
\end{solution}

\begin{solution}{exo:m2:market-making-games:4}
Une largeur égale à la patience, 8 points~: $4e^{-1} = 1.47$ par tour, soit 29.4 sur 20 tours.
\end{solution}

\begin{solution}{exo:m2:market-making-games:5}
0.31. Un quart des sommes est supérieur à 40, donc les acheteurs informés représentent 5~\% des arrivées
($0.2 \times 0.25$)~; les traders non informés n'échangent que 29~\% du temps à cette largeur et achètent
la moitié de cela, soit 11.5~\% des arrivées ($0.8 \times 0.29 / 2$). La largeur filtre les traders non informés
davantage que les informés, si bien qu'un achat est plus souvent informé que ne le laisse voir le tableau, soit 20~\%.
\end{solution}

\begin{solution}{exo:m2:market-making-games:6}
La cotation du teneur de marché est échangée précisément quand quelqu'un pense qu'elle est erronée
à son avantage~; conditionnellement à son exécution, la valeur est moins bonne que ne le pensait le teneur de marché.
Coter sur l'estimation non conditionnelle revient à perdre en moyenne face à ceux qui choisissent d'échanger.
\end{solution}

\begin{solution}{exo:m2:market-making-games:7}
12.5 points, pour un gain de 13.8 par partie~: davantage de traders informés poussent la meilleure largeur vers le haut et le
profit vers le bas (11.0 points et 20.6 avec 10~\% d'informés).
\end{solution}

\begin{solution}{exo:m2:market-making-games:8}
Davantage d'échanges rapportent davantage de fourchette auprès des non-informés, mais un marché serré est aussi échangé
chaque fois que les informés savent qu'il est erroné, et ces échanges coûtent plus que la fourchette
ne rapporte. Sous la meilleure largeur, le profit baisse~; à largeur nulle, le teneur perd 22 par partie dans le
tableau du chapitre.
\end{solution}

\begin{solution}{pb:m2:market-making-games:1}
\textbf{1.} Espérance 35, écart-type 8.03, de 6 à 64.
\textbf{2.} Un jeu compte quatre cartes de chaque valeur~: cinq as ou cinq rois sont impossibles.
\textbf{3.} $e^{-10/8} = 0.29$.
\textbf{4.} 30 à 40.
\textbf{5.} La moitié de la largeur, lorsqu'elle est échangée~: $0.29 \times 5 = 1.43$ par tour, provenant
d'un trader non informé qui arrive.
\textbf{6.} 38.13 après un achat~; 31.87 après une vente.
\textbf{7.} Au premier tour, les cotations sont symétriques autour d'une loi a priori symétrique~; plus tard
le centre a bougé, et un passage indique que la somme se situe probablement à l'intérieur des cotations en cours.
\textbf{8.} Elle monte de 35 à environ 39 après les premiers achats, puis baisse vers 31 à mesure que les ventes
arrivent~; les premiers achats étaient du bruit non informé, ou des achats informés contre un marché encore
centré trop bas, et le teneur ne peut pas savoir lequel.
\textbf{9.} 16.4 par partie à 11.5 points contre 12.0 à 14 points sans mise à jour.
\textbf{10.} Sans apprentissage, son centre reste faux, si bien qu'il doit garder davantage de distance
par rapport à la valeur pour perdre moins face aux informés.
\textbf{11.} 11.5 points, 16.4 par partie.
\textbf{12.} 14.1 à 8 points et 15.1 à 15.
\textbf{13.} 8 points, environ 29 par partie.
\textbf{14.} 11.0 points avec 10~\% d'informés (20.6 par partie), 12.5 avec 30~\% (13.8).
\textbf{15.} $e^{-11.5/8} = 0.24$.
\textbf{16.} À partir des mark-outs de ses exécutions~: de combien la valeur évolue à son détriment après
les échanges, par contrepartie et par taille~; des mouvements défavorables persistants signalent un flux informé.
\textbf{17.} Relever le centre, puisque les acheteurs peuvent savoir que la valeur est plus élevée, et envisager
un décalage des cotations~: un cours vendeur plus haut pour vendre moins en position courte, un cours acheteur plus proche pour racheter.
\textbf{18.} Que le centre est une estimation à réviser à chaque échange, que la largeur est un
prix de la sélection adverse, mis en regard de ce que les non-informés accepteront de payer, et que
apprendre du flux vaut autant qu'une fourchette plus large.
\textbf{19.} Résultat nommé~: \emph{marché de la somme des cartes}~: contre un tableau où un cinquième des
traders est informé, une largeur de 11.5 points maximise le P\&L espéré, soit 16.4 par partie.
\textbf{20.} Le teneur paie les échanges conclus avec des gens qui en savent davantage.
\end{solution}

\begin{solution}{iq:m2:market-making-games:1}
La somme a pour espérance 7, de 2 à 12~; personne n'étant mieux informé, il faut coter autour de 7, par
exemple «~6.5 à 7.5~» pour la taille, plus large si celle-ci est importante au regard du risque
supportable, et se demander qui voudrait l'échanger.
\iqlookfor{le centre, une largeur raisonnée, et qui se trouve de l'autre côté.}
\end{solution}

\begin{solution}{iq:m2:market-making-games:2}
Se demander pourquoi~: l'acheteur en sait peut-être davantage. Relever l'estimation à la hausse d'un montant qui dépend
de la probabilité que l'acheteur soit informé~; déplacer le marché vers le haut, élargir peut-être, et considérer
l'inventaire que je détiens désormais en position courte.
\iqlookfor{espérance conditionnelle et inventaire.}
\end{solution}

\begin{solution}{iq:m2:market-making-games:3}
Glosten et Milgrom~: un teneur de marché neutre au risque, face à des traders informés et non informés, fixe le
cours vendeur à la valeur espérée sachant un achat, et le cours acheteur à la valeur espérée sachant une vente~;
puisque les achats sont plus probables lorsque la valeur est élevée, le cours vendeur se situe au-dessus du cours acheteur même avec
des coûts nuls et un profit espéré nul.
\iqlookfor{espérances conditionnelles selon le sens de l'échange.}
\end{solution}

\begin{solution}{iq:m2:market-making-games:4}
Déplacer le marché vers le haut, à la fois parce que les ventes ont pu être faites à des acheteurs informés et parce que
la position courte est un risque~: relever le cours acheteur davantage que le cours vendeur pour attirer les vendeurs et
décourager les acheteurs, jusqu'à ce que l'inventaire revienne dans les limites.
\iqlookfor{information et inventaire poussent dans le même sens.}
\end{solution}

\begin{solution}{iq:m2:market-making-games:5}
À partir de la relation entre les échanges et les mouvements de prix ultérieurs~: estimer de combien les prix
évoluent après les achats et les ventes (l'impact permanent), ou ajuster un modèle de type Glosten--Milgrom ou PIN
à la suite des achats, des ventes et de la fourchette~; les mark-outs par contrepartie donnent une
mesure directe.
\iqlookfor{impact sur les prix et mark-outs comme indices d'information.}
\end{solution}

\begin{solution}{iq:m2:market-making-games:6}
Un serveur qui détient l'état caché de chaque partie et le carnet d'ordres, avec un protocole de messages pour
les cotations, les transactions et les annulations~; des bots en tant que clients avec la même interface que les joueurs humains~;
journaux horodatés de tout pour le rejeu et le calcul des scores~; salles et scénarios définis comme
des données~; contrôles d'équité sur la latence entre joueurs~; analyses postérieures à la partie du P\&L, de la largeur et
de la sélection adverse.
\iqlookfor{état de référence, même interface pour les bots et les humains, rejeu.}
\end{solution}
